% Gerado por tools/tei2latex.py a partir de vol1-divisao-do-ente.xml. Regenerar, não editar.
% Recorte manual: tirada a matéria preliminar que tools/trecho.py arrastou; \pg{331} reposto.

\pg{331}

%%>>> cap-10

\chapter*{CAPITULO SEGUNDO}
\addcontentsline{toc}{chapter}{CAPITULO SEGUNDO}

\textbf{DIVISÃO DO ENTE} \textsc{Summario}: --- Ente necessario e ente contingente. --- Este infinito e ente finito.

--- Ente simples e ente composto. --- Ente immutavel e ente mutavel.

%%>>> art-12

\section*{Artigo I\\ Ente necessario e ente contingente.}
\addcontentsline{toc}{section}{Artigo I --- Ente necessario e ente contingente.}

54. Divisão do ente. --- Todo o ente --- ou é \emph{acto puro}, --- ou consta de \emph{acto} e de \emph{potencia}. Se é \emph{acto puro}, o ente é \emph{necessario}, \emph{infinito}, \emph{simples}, \emph{immutavel}. Se consta de \emph{acto} e de \emph{potencia}, é \emph{contingente}, \emph{finito}, \emph{composto}, \emph{mutavel}. D'onde a divisão de todo o ente --- em \emph{necessario} e \emph{contingente}, --- em \emph{infinito} e \emph{finito}, --- em \emph{simples} e \emph{composto}, --- em \emph{immutavel} e \emph{mutavel}. --- Esta divisão funda-se na diversa perfeição do ente\footnote{Todo o ente é dotado de \emph{perfeição}. --- \emph{Perfeição é toda a realidade, que serve} para fazer ou completar uma coisa... A perfeição pode ser \emph{parcial} e \emph{total}. \emph{É parcial}, se representa \emph{uma} das realidades, de que o ente é dotado. \emph{É total}. se representa o \emph{conjuncto} das mesmas realidades D' onde se vê que a perfeição não é um attributo especial do ente, mas é uma coisa que constitue, adequada ou inadequadamente, o proprio ente. --- Por isso um ente é e diz-se \emph{perfeito}, isto é, \emph{inteiramente feito}, emquanto possue a perfeição \emph{total}, de modo que não lhe falte nenhuma das realidades. necessarias para o seu complemento. --- A' \emph{perfeição} oppõe-se a \emph{imperfeição}, que denota \emph{ausencia} de perfeição. A imperfeição é \emph{privativa} ou \emph{negativa}. \emph{É privativa}, quando é ausencia de uma perfeição \emph{devida}, e diz-se commummente \emph{privação}; tal é a ausencia de um braço no homem. É \emph{negativa}, quando é a ausencia de uma perfeição \emph{não-devida}, e chama-se \emph{negação}; tal é a ausencia das azas no homem, ou da razão na planta.\par A perfeição \emph{total} divide-se em \emph{absoluta} e \emph{relativa}. A \emph{absoluta} é o conjuncto de \emph{todas} as realidades ou perfeições \emph{possiveis}, sem limitação alguma; tal é a perfeição de \textsc{Deus}, que é o proprio Ser subsistente. A \emph{relativa} é o conjuncto das realidades ou perfeições, que são devidas a um certo genero, ou a uma certa especie de entes; tal é a perfeição de um homem, que tem tudo o que a sua natureza exige. --- A perfeição \emph{absoluta} não admitte graus, nem divisões; a \emph{relativa} admitte graus, --- emquanto um genero ou uma especie pode abranger um maior ou menor numero de perfeições parciaes com relação a um outro genero, ou a uma outra especie, --- ou emquanto um individuo possue uma realidade ou perfeição por um modo mais \emph{intenso} que o\pgn{332} modo por que a possue um outro individuo do mesmo genero, ou da mesma especie. --- A perfeição \emph{total} pode considerar-se --- em relação á \emph{entidade}, quando o ente tem tudo o que é indispensavel para a sua essencia e integridade; --- em relação á \emph{operação}, quando o ente tem tudo o que é necessario para operar em conformidade com a sua essencia; --- em relação ao \emph{fim}, quando o ente, tendo alcançado o fim pelas suas operações, encontra nelle o seu complemento, pois o fim é bem (Cf. \emph{Sum. Th.}, p. I, q. 6, a. 3).\par A perfeição \emph{relativa} divide-se em \emph{simples} e \emph{mista}. --- \emph{Simples} é a que no seu conceito não encerra imperfeição alguma e não se oppõe a perfeições mais elevadas; tal é a \emph{vida}, a \emph{bondade}, a \emph{belleza}, etc. A perfeição \emph{simples} costuma chamar-se tambem \emph{absoluta}, mas emquanto o seu conceito não denota imperfeição. --- \emph{Mista} é a que, embora no seu conceito exprima uma entidade positiva, todavia encerra alguma imperfeição e se oppõe a uma perfeição mais elevada; tal é a \emph{corporeidade}, que, embora denote uma realidade, comtudo importa uma imperfeição, pois o \emph{corpo} é \emph{composto}, \emph{dependente}, etc., --- e se oppõe á \emph{espiritualidade}, que é uma coisa mais elevada do que o corpo (Cf. Ansel., Monol., c. 15).\par A perfeição pode convir ao ente por tres modos: \emph{formal}, \emph{virtual} e \emph{eminente}. Convem \emph{formalmente}, quando existe no ente segundo a propria forma ou razão especifica, isto é, segundo o seu conceito e a sua definição; assim encontra-se o \emph{calor no fogo}, a \emph{sensibilidade} no \emph{cão}, a \emph{razão} no \emph{homem}. Convem \emph{virtualmente}, quando não existe no ente segundo a propria forma, mas existe na virtude ou causalidade do ente, emquanto este é capaz de a produzir; assim encontra-se a \emph{planta} na \emph{semente}, e todo o \emph{effeito} na sua \emph{causa}. Convem \emph{eminentemente}; quando existe no ente, nem segundo a propria forma, nem segundo a energia ou causalidade, mas por um modo mais nobre e mais elevado que o modo por que se encontra nos seres, que a possuem \emph{formalmente}; assim encontram-se em \textsc{Deus} todas as perfeições dos entes creados. --- Note-se que um d'esses modos não exclue outro. Uma perfeição pode convir a um ente por um modo \emph{formal} e \emph{eminente} ao mesmo tempo, e convem quando existe nelle segundo a propria forma, mas num grau mais elevado; assim a \emph{sabedoria}, que é o conhecimento das coisas pelas causas supremas, attribue-se a \textsc{Deus} --- não só \emph{formalmente}, emquanto, em \textsc{Deus}, como no homem, importa o conhecimento pelas causas supremas, --- mas tambem \emph{eminentemente}, emquanto em \textsc{Deus} a \emph{sabedoria} é uma e a mesma coisa com essencia, ao passo que no homem é uma qualidade accidental e limitada. Como tambem uma perfeição pode convir a um ente por um modo \emph{virtual} e \emph{eminente} ao mesmo tempo, e convem quando não existe nelle segundo a propria forma, mas existe na virtude ou capacidade do ente, e de um modo mais elevado e nobre; assim o \emph{raciocinio}, importando um discurso ou uma passagem da potencia para o acto, não se attribue a \textsc{Deus} \emph{formalmente}, --- mas attribue-se \emph{virtualmente}, emquanto a intelligencia divina tem a capacidade de perceber a conclusão no principio, --- e tambem \emph{eminentemente}, emquanto a intelligencia\pgn{333} divina não passa, como a do homem, do conhecido para o desconhecido, mas vê a conclusão no principio sem discurso, por intuição.\par Premissas estas noções, vejamos a razão da divisão do ente. A perfeição pode encontrar-se no ente --- ou de um modo \emph{necessario} ou \emph{contingente}, e temos o ente \emph{necessario} e o ente \emph{contingente}, --- ou de um modo \emph{illimitado} ou \emph{limitado}, e temos o ente \emph{infinito} e o ente \emph{finito}, --- ou de um modo \emph{simples} ou \emph{composto}, e temos o ente \emph{simples} e o ente \emph{composto}, --- ou de um modo \emph{permanente} ou \emph{variavel}, e temos o ente \emph{immutavel} e o ente \emph{mutavel}. --- A primeira especie da divisão refere-se á \emph{existencia} do ente; a segunda, á sua \emph{comprehensão}; a terceira, á sua \emph{constituição}; a quarta, á \emph{permanencia} ou \emph{estabilidade} da perfeição. --- Como se vê, todas estas especies denotam a mesma coisa e importam os mesmos extremos, considerados sob differentes aspectos e significados por varios conceitos e termos, dos quaes um é mais claro e mais obvio que outro: porque o Ente \emph{necessario}, \emph{infinito}, \emph{simples}, \emph{immutavel} é \textsc{Deus}; e o ente \emph{contingente}, \emph{finito}, \emph{composto}, \emph{mutavel} é a creatura. Por isso esta divisão é a mais transcendental que se pode dar. --- Advertimos que, tratando aqui do Ente \emph{necessario}, \emph{infinito}, \emph{simples} e \emph{immutavel}, não entendemos demonstrar a sua \emph{existencia}; mas, admittida, pela experiencia, a existencia de entes \emph{contingentes}, \emph{finitos}, \emph{compostos} e \emph{mutaveis}, procuramos descobrir o \emph{nexo}, que elles têm com o Ente \emph{necessario}, \emph{infinito}, \emph{simples} e \emph{immutavel}.}.
